% GENERATED FILE. Edit canonical lessons, then run npm run generate:books.
% canonical-source-hash: fnv1a64:94bdc60a254aadae
% canonical-lessons: HI-C419-tiraha, HI-C419-cakkar, HI-C419-nisan, HI-C419-imarat, HI-C419-bord

\chapter{Landmarks}
\label{ch:hi-landmarks}
\hlchaptermodality{419}

\begin{chapteropening}
\textbf{By the end of this chapter:} \emph{I can find the way by a T-junction, a mark on a wall, a signboard or a white building.}

Complete the last lesson of chapter 419: I can find the way by a T-junction, a mark on a wall, a signboard or a white building.
\end{chapteropening}

\section[\texorpdfstring{\hi{तिराहा} (tirāhā)}{तिराहा (tirāhā)}]{\hi{तिराहा} (tirāhā) — a T-junction, a meeting of three roads}

\label{lesson:HI-C419-tiraha}



\begin{quote}
Before the new one: say the Hindi for beside, then the Hindi for to deliver.
\end{quote}

\subsection*{You'll want to know: \hi{तिराहा}}
\textbf{\hi{तिराहा}} — \emph{tirāhā} — \textquotedblleft{}a T-junction, a meeting of three roads\textquotedblright{}.

\begin{quote}
\hi{तिराहे} \hi{से} \hi{बाएँ} \hi{मुड़िए।}
\end{quote}

\begin{quote}
\emph{tirāhe se bāẽ muṛiye.}
\end{quote}

\begin{quote}
Turn left at the T-junction.
\end{quote}

Masculine.

\begin{cousinweb}
\textbf{\hi{तिराहा}} (\emph{tirāhā}) is the three-way partner of \hi{चौराहा}: \emph{ti-}, \textquotedblleft{}three\textquotedblright{}, as in \hi{तिकोना}, with \emph{rāh}, \textquotedblleft{}road\textquotedblright{}, usually traced to Persian \emph{rāh}.
\end{cousinweb}

\begin{grammarlens}[title={What you've built}]
One more word for getting around.
\end{grammarlens}

\subsection*{Your turn}
\begin{itemize}
\raggedright
  \item \emph{Say it:} \emph{tirāhā}
  \item \emph{Say it:} \emph{tirāhā}, once more
  \item \emph{Say it:} \emph{tirāhe se bāẽ muṛiye.} — Turn left at the T-junction.
  \item \emph{Recall:} say the Hindi for beside, then the Hindi for to deliver, then say \emph{tirāhā} again
\end{itemize}

\begin{tcolorbox}[breakable,colback=teal!4,colframe=teal!35!black,title={Before you move on}]
What does \hi{तिराहा} mean? (\textquotedblleft{}A T-junction, a meeting of three roads\textquotedblright{}.) Say it once more.
\end{tcolorbox}
\section[\texorpdfstring{\hi{चक्कर} (cakkar)}{चक्कर (cakkar)}]{\hi{चक्कर} (cakkar) — a round; a detour; dizziness}

\label{lesson:HI-C419-cakkar}



\begin{quote}
Before the new one: say the Hindi for to deliver, then the Hindi for a T-junction.
\end{quote}

\subsection*{You'll want to know: \hi{चक्कर}}
\textbf{\hi{चक्कर}} — \emph{cakkar} — \textquotedblleft{}a round; a detour; dizziness\textquotedblright{}.

\begin{quote}
\hi{ग़लत} \hi{रास्ते} \hi{की} \hi{वजह} \hi{से} \hi{हमें} \hi{लंबा} \hi{चक्कर} \hi{लगाना} \hi{पड़ा।}
\end{quote}

\begin{quote}
\emph{ġalat rāste kī vajah se hamẽ lambā cakkar lagānā paṛā.}
\end{quote}

\begin{quote}
Because of the wrong road we had to go a long way round.
\end{quote}

Masculine.

\hi{चक्कर} \hi{आना}, to feel dizzy: \hi{मुझे} \hi{चक्कर} \hi{आ} \hi{रहा} \hi{है}.

\begin{cousinweb}
\textbf{\hi{चक्कर}} (\emph{cakkar}) is a semi-learned form of Sanskrit \emph{cakra}, \textquotedblleft{}a wheel, a circle\textquotedblright{} (CDIAL 4538); the fully inherited Hindi form in that entry is \emph{cāk}, \textquotedblleft{}a wheel, a millstone\textquotedblright{}. A \hi{चक्कर} is any going round: a detour, a round of the market, a spinning head.
\end{cousinweb}

\begin{grammarlens}[title={What you've built}]
One more word for getting around.
\end{grammarlens}

\subsection*{Your turn}
\begin{itemize}
\raggedright
  \item \emph{Say it:} \emph{cakkar}
  \item \emph{Say it:} \emph{cakkar}, once more
  \item \emph{Say it:} \emph{ġalat rāste kī vajah se hamẽ lambā cakkar lagānā paṛā.} — Because of the wrong road we had to go a long way round.
  \item \emph{Recall:} say the Hindi for to deliver, then the Hindi for a T-junction, then say \emph{cakkar} again
\end{itemize}

\begin{tcolorbox}[breakable,colback=teal!4,colframe=teal!35!black,title={Before you move on}]
What does \hi{चक्कर} mean? (\textquotedblleft{}A round; a detour; dizziness\textquotedblright{}.) Say it once more.
\end{tcolorbox}
\section[\texorpdfstring{\hi{निशान} (niśān)}{निशान (niśān)}]{\hi{निशान} (niśān) — a mark left behind; a sign; a scar}

\label{lesson:HI-C419-nisan}



\begin{quote}
Before the new one: say the Hindi for a T-junction, then the Hindi for a round.
\end{quote}

\subsection*{You'll want to know: \hi{निशान}}
\textbf{\hi{निशान}} — \emph{niśān} — \textquotedblleft{}a mark left behind; a sign; a scar\textquotedblright{}.

\begin{quote}
\hi{दीवार} \hi{पर} \hi{लाल} \hi{निशान} \hi{है}, \hi{वहाँ} \hi{से} \hi{दाएँ} \hi{मुड़िए।}
\end{quote}

\begin{quote}
\emph{dīvār par lāl niśān hai, vahā̃ se dāẽ muṛiye.}
\end{quote}

\begin{quote}
There is a red mark on the wall; turn right there.
\end{quote}

Masculine.

\begin{cousinweb}
\textbf{\hi{निशान}} (\emph{niśān}) is usually traced to Persian \emph{nišān}, \textquotedblleft{}a sign, a mark\textquotedblright{}: a scar, a stain, a mark to steer by.
\end{cousinweb}

\begin{grammarlens}[title={What you've built}]
One more word for getting around.
\end{grammarlens}

\subsection*{Your turn}
\begin{itemize}
\raggedright
  \item \emph{Say it:} \emph{niśān}
  \item \emph{Say it:} \emph{niśān}, once more
  \item \emph{Say it:} \emph{dīvār par lāl niśān hai, vahā̃ se dāẽ muṛiye.} — There is a red mark on the wall; turn right there.
  \item \emph{Recall:} say the Hindi for a T-junction, then the Hindi for a round, then say \emph{niśān} again
\end{itemize}

\begin{tcolorbox}[breakable,colback=teal!4,colframe=teal!35!black,title={Before you move on}]
What does \hi{निशान} mean? (\textquotedblleft{}A mark left behind; a sign; a scar\textquotedblright{}.) Say it once more.
\end{tcolorbox}
\section[\texorpdfstring{\hi{इमारत} (imārat)}{इमारत (imārat)}]{\hi{इमारत} (imārat) — a building}

\label{lesson:HI-C419-imarat}



\begin{quote}
Before the new one: say the Hindi for a round, then the Hindi for a mark left behind.
\end{quote}

\subsection*{You'll want to know: \hi{इमारत}}
\textbf{\hi{इमारत}} — \emph{imārat} — \textquotedblleft{}a building\textquotedblright{}.

\begin{quote}
\hi{वह} \hi{सफ़ेद} \hi{इमारत} \hi{अस्पताल} \hi{है।}
\end{quote}

\begin{quote}
\emph{vah safed imārat aspatāl hai.}
\end{quote}

\begin{quote}
That white building is the hospital.
\end{quote}

Feminine.

\begin{cousinweb}
\textbf{\hi{इमारत}} (\emph{imārat}) is usually traced to Arabic \emph{imāra}, \textquotedblleft{}a building\textquotedblright{}, through Persian, which gave it its final t; it is feminine.
\end{cousinweb}

\begin{grammarlens}[title={What you've built}]
One more word for getting around.
\end{grammarlens}

\subsection*{Your turn}
\begin{itemize}
\raggedright
  \item \emph{Say it:} \emph{imārat}
  \item \emph{Say it:} \emph{imārat}, once more
  \item \emph{Say it:} \emph{vah safed imārat aspatāl hai.} — That white building is the hospital.
  \item \emph{Recall:} say the Hindi for a round, then the Hindi for a mark left behind, then say \emph{imārat} again
\end{itemize}

\begin{tcolorbox}[breakable,colback=teal!4,colframe=teal!35!black,title={Before you move on}]
What does \hi{इमारत} mean? (\textquotedblleft{}A building\textquotedblright{}.) Say it once more.
\end{tcolorbox}
\section[\texorpdfstring{\hi{बोर्ड} (borḍ)}{बोर्ड (borḍ)}]{\hi{बोर्ड} (borḍ) — a signboard; a board}

\label{lesson:HI-C419-bord}



\begin{quote}
Before the new one: say the Hindi for a mark left behind, then the Hindi for a building.
\end{quote}

\subsection*{You'll want to know: \hi{बोर्ड}}
\textbf{\hi{बोर्ड}} — \emph{borḍ} — \textquotedblleft{}a signboard; a board\textquotedblright{}.

\begin{quote}
\hi{बोर्ड} \hi{पर} \hi{लिखा} \hi{है} \hi{कि} \hi{अस्पताल} \hi{दाईं} \hi{तरफ़} \hi{है।}
\end{quote}

\begin{quote}
\emph{borḍ par likhā hai ki aspatāl dāīṁ taraf hai.}
\end{quote}

\begin{quote}
The board says the hospital is on the right.
\end{quote}

Masculine.

\begin{cousinweb}
\textbf{\hi{बोर्ड}} (\emph{borḍ}) is borrowed from English \emph{board}, above all the signboard over a shop or at a turning.
\end{cousinweb}

\begin{grammarlens}[title={What you've built}]
One more word for getting around.
\end{grammarlens}

\subsection*{Your turn}
\begin{itemize}
\raggedright
  \item \emph{Say it:} \emph{borḍ}
  \item \emph{Say it:} \emph{borḍ}, once more
  \item \emph{Say it:} \emph{borḍ par likhā hai ki aspatāl dāīṁ taraf hai.} — The board says the hospital is on the right.
  \item \emph{Recall:} say the Hindi for a mark left behind, then the Hindi for a building, then say \emph{borḍ} again
\end{itemize}

\begin{tcolorbox}[breakable,colback=teal!4,colframe=teal!35!black,title={Before you move on}]
What does \hi{बोर्ड} mean? (\textquotedblleft{}A signboard; a board\textquotedblright{}.) Say it once more.
\end{tcolorbox}
